% sections/00_abstract.tex -- 摘要

本报告基于黄仁勋（NVIDIA CEO）与埃隆·马斯克（SpaceX CEO）在2024--2026年间的12条一手公开言论，从科学物理、工程可行性与商业经济三个维度对太空（轨道）算力进行了独立技术研判。报告覆盖11个关键子问题（Q1--Q11），采用斯特藩-玻尔兹曼定律定量计算、PUE对比分析、LEO热环境建模、TCO全生命周期对比、GPU能效历史外推与多变量敏感性分析等方法。

核心发现如下：

\textbf{物理层面}：SpaceX AI1卫星宣称的1,400\unitWpm{}散热密度在物理上成立（双面辐射、面板温度约342.5K、辐射率$\varepsilon\approx0.90$），但已逼近当前材料天花板。1~GW轨道数据中心需要约0.71~km\textsuperscript{2}双面散热面积（约6,400颗AI1级卫星），散热面积需求巨大但非物理不可能。LEO真实热环境（太阳辐照、地球反照、红外辐射、90分钟昼夜循环）使得散热器实际效率比理想真空计算低15--40\%。

\textbf{工程层面}：Starship V3可将发射成本降至\$100--200/kg（基准），部署1MW轨道算力约需1次发射、\$8--18M发射成本。但轨道硬件面临严重可靠性挑战——消费级GPU无法在轨道存活，NVIDIA Rubin Space-1虽具备辐射耐受设计但长期可靠性未知。在轨GPU级维修目前完全不可行。

\textbf{商业层面}：2026年轨道数据中心10年TCO约为地面的10.5倍（\$681M/MW vs \$37.8--91.4M/MW，中位\$64.6M）。基准情景TCO \$681M/MW（发射\$200/kg、硬件溢价3.5x、寿命5年）；乐观情景降至\$261.4M/MW（4.3x）；悲观情景\$1,220M/MW（17.7x）。发射成本即使降至\$50/kg，轨道TCO仍为\$678.6M（10.5倍）——单靠发射成本下降不足以翻转。翻转条件组合（发射$<$\$100/kg + 硬件溢价$<$3x + 寿命$>$7年）下，轨道TCO仍为\$408.1M（6.3x），未达地面水平。乐观估计2029--2032年可能达到临界可行（综合概率$<$25\%）。

\textbf{综合判断}：当前太空算力在物理上可能、在工程上脆弱、在商业上不可行。黄仁勋的"务实保留"比马斯克的"5年内更便宜"在工程层面更站得住脚，但马斯克的乐观假设在技术曲线上并非完全不可能。核心不确定性不在于任何单一参数，而在于多参数同步改善的概率。
